\scanpage{213}
\hfill -- \mspage{49} -- \hfill IX

On peut évidemment supposer $S$ strictement local. Avec les notations de 4.1 et de 4.2, prenons $K'$ le plus petit possible, i.e. supposons que $\Gamma$ opère \uline{fidèlement} sur $A_o^{*\unsure{0|o}}$. Par hypothèse, $\Gamma$ opère trivialement sur ${}_nA(\overline K)$, et à fortiori sur ${}_nA(\overline K)^{ef}\simeq{}_nA_o^*(k)\supset{}_nA_o^{*\unsure{0|o}}(k)$. On conclut alors par les deux lemmes suivants :

\uline{Lemme} 4.7.1. (J.P. SERRE). \uline{Soient} $G$ \uline{un schéma en groupes commutatif sur un corps} $k$, \uline{extension d'un schéma abélien par un tore}, $n$ \uline{un nombre entier} $\geq2$. \uline{Alors tout automorphisme d'ordre fini de} $G$ \uline{induisant l'identité sur} ${}_nG$ \uline{est réduit à l'identité si} $n\neq2$, \uline{et est l'identité ou d'ordre deux si} $n=2$.

\uline{Lemme} 4.7.2. \uline{Un groupe fini dont tout élément} $\neq1$ \uline{est d'ordre} $2$ \uline{est isomorphe à} $(\mathbb Z/2\mathbb Z)^r$, \uline{pour un} $r\geq0$ \uline{convenable}.

4.7.3. Rappelons la démonstration de 4.7.1, qui est bien connue. On note que grâce au fait que $n\operatorname{id}_G$ est un épimorphisme, l'hypothèse que $u$ induise l'identité sur ${}_nG$ i.e. que $1-u$ induise $0$ sur ${}_nG$, s'exprime par la condition que l'on ait $1-u=nv$ dans l'anneau des endomorphismes de $G$. Pour prouver la conclusion annoncée, on est ramené (en décomposant le groupe cyclique fini engendré par $u$ en ses composantes $p$-primaires, pour les divers nombres premiers $p$ divisant l'ordre de $u$) au cas où $u$ est d'ordre $p^N$, puissance d'un nombre premier. Soit $E$ le sous-anneau de $\operatorname{End}(G)$ engendré par $v$, qui est un quotient de $\mathbb Z[T]/(T^{p^N}-1)$, donc de type fini sur $\mathbb Z$, et qui de plus est sans torsion (ceci est conséquence formelle du fait que $G$ est divisible). Il se plonge donc dans $E_{\mathbb Q}=E\otimes_{\mathbb Z}\mathbb Q$, qui
\margin{367}
\margin{367}
\margin{212}
